%% feature/sec-abstract -- body only; wrapped in rgabstractbox in main.tex
Reasoning-capable language models often improve final-answer accuracy by generating long chains of thought, but this improvement comes with substantial inference cost. Two recent findings suggest that the form and repetition of context strongly affect model behavior. Prompt repetition, duplicating the user query at least once, improves non-reasoning language models, while distilled reasoning content placed before a problem guides reasoning models: I2S shows that abstract insights help where raw reasoning traces hurt, and TRS shows that compact, reusable skill cards preserve accuracy while substantially reducing generated tokens. This paper studies the unexplored intersection of these findings: if a fixed, already-useful insight is available for a problem, does repeating the same insight multiple times improve, saturate, or hurt performance and token efficiency?

We propose a controlled evaluation that isolates this repetition effect. Each problem is drawn from the released TRS Reasoning-Skill dataset, which provides a target question, a gold answer, and a precomputed skill card that is independently validated to help at a single insertion. We repeat the identical card $k \in \{0,1,2,3,5\}$ times before the problem, holding the question, insight text, prompt template, model, and decoding settings fixed, so the repetition count is the only changing variable and the result cannot be confounded with retrieval or extraction quality. We compare against direct prompting, pure prompt repetition, and raw reasoning-trace baselines at matched copy counts, and evaluate several reasoning-capable model families through a unified API with full input, completion, and reasoning-token accounting. Because prior work shows distilled insights primarily affect tokens rather than accuracy, token efficiency is our primary axis of interest. --- Results ---
